\section{Discussion}
\label{sec:discussion}

\subsection{From Single-Day to Multi-Day: A Unified Interpretation}

The progression from single-day validation to multi-day regime detection reveals a coherent picture of LLM structural reasoning capabilities. Single-day analysis establishes that LLMs detect dealer constraints through genuine mechanism identification: the raw chain superiority (92.3\% vs.\ 61.5\%) proves the model reconstructs dealer positioning from first principles, while the inverse p-hacking finding (33\% lower range expansion) and alpha orthogonality (Sharpe 1.8 $\rightarrow$ 0.1) confirm the model identifies structural states rather than exploitable anomalies.

These single-day findings provide necessary grounding for interpreting regime detection results. The 69.1pp discrimination between 2024 and 2020 is credible precisely because we have established that the LLM performs genuine structural analysis---not pattern matching on pre-calculated metrics. The alpha disappearance puzzle from single-day work motivates the regime extension: if dealers maintain persistent positioning, the structural constraint is detectable but not tradeable because it never transitions. Multi-day analysis confirms this interpretation: 100\% detection in 2024--2025 reflects sustained structural coherence rather than detection-profitable switching.

Together, both temporal scales support the same conclusion: LLMs identify emergent constraints in complex adaptive systems---coordination patterns arising from decentralized dealer inventory management~\cite{grossman1988liquidity}---without capturing the dynamic adaptations that generate trading alpha. This positions temporal obfuscation testing as a domain-agnostic methodology for validating structural reasoning wherever training data contamination threatens evaluation validity.

\subsection{LLM Temporal Reasoning at 30-Day Scales}

The shift from 5-day to 30-day analysis windows fundamentally transforms the detection task. Earlier 5-day trajectory testing achieved 98--100\% detection across all market conditions, identifying daily hedging flows---mechanics present in all markets since Black-Scholes~\cite{black1973pricing}---rather than persistent regimes. The move to 30-day windows with strict persistence criteria ($\geq$70\% same-sign dominance, $\geq$\$5B magnitude, $\leq$5 flips) converted the task from universal pattern matching to selective regime identification. This selectivity is precisely the methodological advance: detection rates ranging from 12\% to 81\% demonstrate discrimination rather than trivial universal recognition.

Critically, LLM reasoning quality remained high across both detection and non-detection conditions. Manual review of 50 randomly sampled responses (25 from 2024, 25 from 2020) revealed 98\% mechanical accuracy, 100\% criterion citation, and 88\% step-by-step calculation verification. The LLM accurately computed persistence percentages, magnitude averages, and flip counts regardless of whether the window qualified as a persistent regime. When a window fails regime criteria, the model correctly identifies \textit{why} it fails---insufficient magnitude, excessive sign flips, or borderline persistence---demonstrating criterion-based evaluation rather than binary classification heuristics.

This reasoning quality finding has important implications for interpretability. Unlike black-box classifiers that output binary labels, the LLM provides transparent justifications that can be audited by domain experts. Each response traces the logical chain from raw GEX values through persistence calculation to regime determination, producing an explanation that is both mechanically verifiable and economically interpretable. The 98\% mechanical accuracy across both detection conditions suggests the model reasons structurally about dealer constraints rather than pattern matching on summary statistics.

\subsection{Regime Exceptionality vs.\ Universal Positioning}

A critical methodological concern is whether our framework detects \textit{exceptional persistent regimes} or simply identifies universal dealer positioning present in all markets. Four empirical findings establish regime exceptionality:

\textbf{Cross-period discrimination}: The 69.1 percentage point difference between 2024 detection (81.2\%) and 2020 detection (12.1\%) demonstrates that the framework identifies selective structural persistence. If dealer gamma positioning were always regime-like, detection rates would converge across periods rather than exhibiting 5.7$\times$ separation ($p < 0.0001$, Cram\'{e}r's $\phi = 0.672$). The effect size exceeds conventional thresholds for ``large'' effects ($\phi > 0.5$), indicating a substantive structural difference rather than a marginal statistical artifact.

\textbf{2020 as negative baseline}: The 12.1\% detection rate in 2020 establishes a crucial baseline---dealer gamma hedging was active (average \$2.85B magnitude), but 87.9\% of windows failed persistence criteria. This indicates the framework correctly rejects fragmented positioning as non-regimes, even when dealer constraints are present. 2020 serves as an implicit negative control: active hedging dynamics without sustained regime structure.

\textbf{Magnitude threshold enforcement}: The \$5B average magnitude criterion excludes routine dealer activity. In 2020, 83.3\% same-sign persistence existed but at \$2.85B scale---insufficient to create regime-level market impact. By requiring both persistence \textit{and} magnitude, the framework filters low-conviction positioning that does not materially constrain price dynamics. This dual requirement prevents the framework from detecting the ubiquitous daily hedging flows that earlier 5-day testing identified universally.

\textbf{Deliberate methodological pivot}: We deliberately shifted from 5-day windows (98--100\% universal detection) to 30-day windows with stricter criteria to transform the task from universal pattern detection to exceptional regime identification. The resulting 12--81\% detection range validates this pivot. The framework discriminates sustained structural persistence (2024) from fragmented positioning (2020), confirmed through comprehensive negative controls (0\% false positive on transitional and low-magnitude synthetics) and cross-period comparison.

Collectively, these findings address the ``always a regime'' objection by demonstrating that our framework enforces criteria rather than merely detecting ``negative GEX.'' The regime sign flip between periods---positive dominant in 2020, negative dominant in 2024---further demonstrates sensitivity to directional structure, not just magnitude.

\subsection{Alternative Explanations and Robustness}

\subsubsection{Price Normalization and Asset Inflation}

A critical methodological concern is whether the 2020$\rightarrow$2024 detection increase (12.1\% $\rightarrow$ 81.2\%) reflects genuine structural change or merely price inflation artifacts. SPY appreciated approximately 2.3$\times$ over this period (330 $\rightarrow$ 500), and GEX calculations scale with underlying price squared. To distinguish structural shifts from asset inflation, we applied the SPY price growth factor $\left(\frac{500}{330}\right)^2 \approx 2.30\times$ to all 2020 GEX magnitudes, simulating what detection rates would have been if 2020 traded at 2024 price levels.

Price normalization increased 2020 detection from 12.1\% to 33.2\% (+21.1 percentage points), accounting for approximately 31\% of the total 69.1pp gap. The remaining 48.0pp (69\%) persists after normalization, attributable to structural evolution in dealer positioning. Criterion-level analysis reveals the structural nature of this gap across all three regime dimensions:

\begin{itemize}
    \item \textbf{Persistence}: 96\% (2024) vs.\ 69\% (normalized 2020), +27pp---persistence gains cannot be explained by inflation adjustment alone, as sign dominance is scale-invariant.
    \item \textbf{Magnitude}: 81\% (2024) vs.\ 33\% (normalized 2020), +48pp---even after the 2.3$\times$ handicap, 2024 remains sharply higher, indicating real growth in dealer positioning scale.
    \item \textbf{Stability}: 99\% (2024) vs.\ 54\% (normalized 2020), +45pp---directional consistency is structurally different, with 2024 exhibiting nearly zero sign flips per window compared to 2020's frequent reversals.
\end{itemize}

Even after giving 2020 a 2.3$\times$ magnitude handicap, it still fails persistence and stability criteria at far higher rates than 2024. The 2.2:1 structural-to-price ratio confirms that the 2023$\rightarrow$2024 transition reflects fundamental market restructuring, not threshold artifacts from asset appreciation. This analysis demonstrates transparency regarding the inflation confound while strengthening the structural shift thesis.

\subsubsection{Market Structure Evolution: The 0DTE Era}

The 69.1 percentage point detection difference coincides temporally with 0DTE options proliferation (2020 $<$5\% volume $\rightarrow$ 2024 $\approx$48\%)~\cite{cboe2024zero}. Academic research confirms these structural effects: Brogaard et al.~\cite{brogaard2023does} find that a one standard deviation increase in 0DTE trading leads to a 10.4\% increase in market volatility. While this temporal alignment provides strong circumstantial evidence, observational data cannot strictly prove causality without a counterfactual control. We must systematically address major confounders that changed simultaneously from 2020 to 2024.

\textbf{Confounder 1: Interest Rates} (Federal Funds: 0.25\% $\rightarrow$ 5.5\%). Higher rates increase funding costs for multi-day hedges and may drive dealers toward shorter-duration instruments including 0DTE contracts. Rising rates also increase option Rho sensitivity, making long-dated options more expensive to hedge. However, interest rate changes are a secular monotonic factor (continuously rising 2020--2023), whereas our detection pattern is non-monotonic (3.7\% $\rightarrow$ 32.4\% $\rightarrow$ 20.2\% $\rightarrow$ 100\%). The 2023$\rightarrow$2024 regime transition coincides with the Federal Reserve's terminal rate plateau in December 2023, not ongoing increases. This temporal mismatch argues against interest rates as the primary structural driver, though they likely contributed as a secondary amplifying factor.

\textbf{Confounder 2: Volatility Regimes.} 2020 included COVID-19 crisis volatility (VIX spike to 80+), while 2024 exhibited subdued volatility (VIX mostly 12--20). Regime persistence might improve naturally in low-volatility environments independent of 0DTE. However, if volatility regime were the primary driver, we would expect higher detection in the volatile 2023 environment (post-FOMC uncertainty, regional banking stress, debt ceiling concerns) and lower detection in calmer 2024---the opposite of what we observe. The 2023 detection dip (32.4\% $\rightarrow$ 20.2\%) followed by the 2024 explosion (100\%) suggests a tipping point dynamic: regime formation requires both sustained hedging pressure \textit{and} a volatility environment permitting consolidation. The non-monotonic pattern strengthens rather than weakens the 0DTE hypothesis, as it demonstrates that volume alone is insufficient---only when 0DTE volumes reached saturation ($\approx$48\%) \textit{and} volatility subsided did positioning constraints become binding across all market conditions.

\textbf{Confounder 3: Quantitative Trading Growth.} Secular growth in algorithmic market making and increased options liquidity have been accelerating since 2015. If positioning persistence improved gradually from technological evolution, we would expect steady detection rate progression across the sample period. Instead, detection remains flat through 2023 (3.7\% $\rightarrow$ 32.4\% $\rightarrow$ 20.2\%) then exhibits a discontinuous jump to 100\% in 2024. This cliff-like transition cannot be explained by gradual technological forces. Market maker concentration---with firms such as Citadel Securities and Virtu Financial dominating equity options~\cite{cboe2024market}---may amplify 0DTE effects by producing more coordinated residual positioning than a fragmented dealer population. We acknowledge concentration as a legitimate contributing factor while noting the temporal mismatch argues against it as the primary driver: industry consolidation has been gradual since 2015, whereas the detection pattern shows a sharp 2023$\rightarrow$2024 transition.

The discontinuous detection pattern, temporal alignment with 0DTE market saturation, and independent OCC volume corroboration~\cite{occ2022record} together provide strong circumstantial evidence that 0DTE proliferation is the most plausible primary catalyst for the structural regime transition. Interest rates, volatility regimes, and quantitative trading growth are secondary factors that may have amplified the transition but cannot individually explain its timing or magnitude.

\subsubsection{The Shuffle Test Boundary Condition}

The Phase 2a shuffle test warrants careful interpretation. The 61.1\% false positive rate on 2024 shuffled data appears concerning but is itself an important structural finding: 2024 regimes exhibit such extreme same-sign dominance (96\%) that randomizing day order rarely breaks the persistence threshold. This validates that 2024 detection reflects aggregate statistical properties robust to reordering, while 2020 detection (12.1\% on shuffled data) depends on temporal structure---consistent with fragmented markets where day ordering matters.

The shuffle test thus identifies a known boundary condition: the persistence criterion alone is insufficient for highly polarized markets where structural dominance survives randomization. This is documented as a limitation rather than a framework failure. Future work should develop supplementary criteria---such as autocorrelation structure or within-window trend analysis---that distinguish genuine temporal regimes from statistically dominant distributions even in extreme environments.

\subsubsection{Formula Agreement Testing}

To validate that detection is robust to mathematical formulation, we executed a Formula Agreement Test on the Q1 2024 validation set (61 windows), comparing absolute dollar-scaled GEX ($S^2$ formulation) against normalized GEX ratios (values scaled between $-1.0$ and $+1.0$, stripping absolute dollar scale). The test achieved 100\% agreement (61/61 windows): every window detected by the dollar-scaled model was also detected by the normalized model, and every rejected window was mutually rejected.

This confirms that the LLM identifies regimes through structural patterns---persistence, sign consistency, and flip frequency---which remain robust across mathematical transformations. While absolute magnitude provides helpful ``economic significance'' cues for human analysts, the model's regime classification depends on structural properties rather than nominal scale. Combined with the price normalization analysis, these results establish that \textit{structure} (persistence and stability) is the dominant variable discriminating regime states, with dollar magnitude serving as a proxy for structural intensity rather than the causal factor. This resolves a potential ``magnitude paradox'': 2020 fails to qualify as a regime even when adjusted for inflation (lacking structural persistence), and 2024 qualifies even when stripped of dollar magnitude (possessing robust structure).

\subsubsection{Threshold Sensitivity}

Threshold sensitivity analysis (Section~\ref{sec:regime}) provides the strongest robustness evidence: re-classifying all 1,412 windows under 29 threshold configurations yields 100\% detection for 2024 across every variation, with effect sizes ranging from +68.5pp to +100.0pp. This confirms that the $\geq$70\% persistence, $\geq$\$5B magnitude, and $\leq$5 flip thresholds are not cherry-picked---the post-0DTE regime is so structurally dominant that no reasonable parameter choice alters classification.

\subsection{The Alpha Disappearance Puzzle}

A central finding of this study---stable detection capability (68--74\%) persisting while economic profitability collapses (Sharpe 1.8 $\rightarrow$ 0.1, Fig.~\ref{fig:quarterly})---raises an important question we deliberately leave for future investigation: \textit{why did alpha disappear?}

We document the phenomenon but do not claim to explain it. Candidate explanations include:

\begin{enumerate}
    \item \textbf{Market adaptation}: As participants recognize and arbitrage the persistent negative gamma regime, the pricing inefficiency that generated alpha in early quarters is gradually eliminated. The structural constraint remains (dealers are still short gamma), but its market impact is anticipated and priced in.
    \item \textbf{Crowding effects}: As more traders exploit similar structural signals---whether through LLM analysis, gamma exposure services~\cite{spotgamma2021}, or dealer positioning analytics~\cite{jpmorgan2023}---the resulting crowded positioning compresses returns toward zero.
    \item \textbf{0DTE microstructure changes}: Daily 0DTE volumes grew throughout 2024, potentially altering the intraday hedging dynamics that generated exploitable patterns in earlier quarters. The structural constraint intensified but its expression shifted from predictable daily patterns to more diffuse intraday dynamics.
    \item \textbf{Regime-specific dynamics}: Persistent states offer detection opportunity but no transition-based trading edge. Alpha may require detecting regime \textit{transitions} rather than static structural states---the persistent negative gamma environment provides no informational asymmetry because all participants experience the same positioning pressure continuously.
\end{enumerate}

The fourth explanation connects directly to the regime detection results. The persistent negative gamma environment of 2024 represents an extreme case where the constraint mechanism is detectable but not tradeable precisely because it never switches off. Multi-day analysis supports this interpretation: 100\% detection in 2024--2025 reflects a market where structural constraints are so pervasive that they provide no informational edge. This detection-profitability divergence has broad implications for any practitioner deploying structural pattern detection in trading systems---genuine understanding of market mechanics does not automatically translate to economic value.

This puzzle motivates future research examining whether alpha generation requires detecting regime transitions rather than static states, and whether the structural intelligence provided by LLM detection can be paired with transition-sensitive statistical models to recover economic value.

\subsection{Practitioner Implications}

For practitioners deploying LLM-based market analysis, our findings suggest specific considerations:

\textbf{Role separation}: LLMs excel at mechanism identification (90.9\% accuracy) but underperform at probability calibration (AUC 0.488 vs.\ statistical model's 0.681). Optimal deployment uses LLMs for structural detection with statistical models providing probability estimates. This hybrid approach leverages each tool's comparative advantage---qualitative reasoning from LLMs, quantitative precision from parametric models.

\textbf{Raw data advantage}: The 30.8pp improvement from raw chain data suggests providing LLMs with granular strike-level data rather than summary statistics. Net GEX may actually \textit{obscure} structural signals visible in the full options surface~\cite{hayek1945}. This has direct implications for data pipeline design: production systems should preserve strike-level granularity rather than pre-aggregating into scalar metrics.

\textbf{Detection does not equal alpha}: The stable detection/declining profitability pattern (Fig.~\ref{fig:quarterly}) demonstrates that structural pattern detection does not guarantee trading edge. Detected patterns serve as risk management inputs---structural constraints that suppress or amplify volatility are precisely the regime information risk models require, even when they provide no directional trading signal.

\textbf{Prompt design}: The 28.5pp biased-unbiased gap reveals significant framing sensitivity. Practitioners should employ minimal prompting for core reasoning, avoiding keyword enumeration that triggers template responses. Our analysis showed that explicitly requesting qualitative language (e.g., ``amplification,'' ``cascading,'' ``dampening'') produced nearly identical template-like responses across high-alpha and zero-alpha quarters, masking genuine contextual adaptation visible in brief, unprompted outputs. Rich structured outputs should be reserved for post-detection explanation phases.

\textbf{Signal quality assessment}: Non-detection day analysis reveals that LLMs require concentrated, coherent gamma signals for reliable detection. Practitioners should monitor GEX concentration metrics alongside detection confidence---fragmented gamma distributions indicate reduced signal reliability. Days where the LLM fails to detect patterns despite negative GEX tend to exhibit dispersed positioning across strikes rather than concentrated gamma at key levels, suggesting that signal coherence matters as much as signal magnitude.

\textbf{Obfuscation testing for validation}: Before deploying LLM-based analysis in production, practitioners should validate model understanding using obfuscation testing---stripping temporal context to confirm structural reasoning rather than memorization. This is particularly critical when applying models to novel market regimes or asset classes not well-represented in training data. If detection rates remain stable under obfuscation, the model is reasoning from structure; if they collapse, the model depends on contextual cues that may not generalize to new environments.

\subsection{Limitations}

Several limitations warrant explicit consideration, though we argue each also clarifies specific aspects of our methodology:

\textbf{Single asset focus}: Testing exclusively on SPY limits generalizability, though SPY provides methodological advantages: it concentrates dealer hedging activity in the world's most liquid options market, exhibits minimal bid-ask spread noise, and serves as the primary 0DTE vehicle (85\%+ of same-day expiration volume)~\cite{cboe2024zero}. Detecting patterns in the most informationally efficient market provides conservative validation that should generalize to less efficient underlyings. However, cross-asset validation---particularly on equities with positive gamma regimes absent from our sample---remains essential future work.

\textbf{Single regime environment}: All 242 single-day test days and 100\% of 2024 regime windows exhibited negative gamma, precluding regime-transition analysis. We frame this as a harder test---detecting mechanisms within uniformity rather than exploiting variation---but it leaves open whether detection capability extends to positive gamma regimes or regime-switching dynamics~\cite{ni2005stock,garleanu2009demand}. The 0DTE-driven structural shift may represent a new market normal, making single-regime validation increasingly relevant, but historical comparison across regime types remains valuable for establishing generalizability.

\textbf{End-of-day granularity}: Daily snapshots miss intraday microstructure dynamics critical for 0DTE hedging, where contracts expire and hedges unwind within hours. However, EOD data achieves AUC = 0.681 for T+1 outcomes ($p = 0.0075$), validating sufficient predictive information for overnight positioning decisions. The statistical significance of this result confirms that end-of-day open interest captures meaningful structural information about dealer positioning, even though it represents only a snapshot of the continuous intraday process. Intraday analysis could improve precision but is not necessary for validating the core detection methodology.

\textbf{Single model family}: Validation using o3-mini and o4-mini cannot distinguish architecture-specific capabilities from general LLM reasoning. These models share the same OpenAI reasoning architecture, and cross-model validation using alternative architectures (e.g., Claude, Gemini, open-source alternatives) would strengthen generalizability claims. However, our results establish that at least one LLM architecture detects structural patterns under obfuscation---a necessary first step before broader cross-architecture validation.

\textbf{Small raw chain sample}: The raw chain validation was conducted on a limited sample of $n = 13$ days, constrained by the availability of complete, high-fidelity raw options chain data. While statistically small, the 30.8pp effect size substantially exceeds the minimum detectable effect for this sample ($\sim$15pp at 80\% power), providing statistical confidence that the raw data advantage is genuine rather than a sampling artifact. Larger-scale replication remains important for precise effect size estimation and for understanding which specific strike-level features drive the improvement.

\textbf{Prompt sensitivity}: The 28.5pp gap between biased and unbiased detection reveals significant framing effects. Rather than undermining our results, this finding motivates the unbiased methodology we employ and establishes prompt design as a critical variable for financial LLM deployment. Researchers designing LLM-based financial analysis should prioritize minimal prompting to preserve reasoning fidelity, as keyword-rich prompts risk eliciting trained templates rather than genuine structural reasoning.

\textbf{Shuffle test boundary condition}: The 61\% false positive rate on 2024 shuffled data reflects extreme regime dominance rather than framework failure, but indicates the persistence criterion alone is insufficient for highly polarized markets where structural dominance survives randomization. This is documented as a known boundary condition of the current methodology. Future work should develop supplementary criteria---such as autocorrelation structure, within-window trend analysis, or intraday dynamics---that distinguish genuine temporal regimes from statistically dominant distributions.

\textbf{Data gaps}: Alpha Vantage provides 4 AM -- 8 PM ET data only, precluding overnight transmission analysis to global sessions (Tokyo 7 PM -- 2 AM ET, London 3 AM -- 4 AM ET). Intraday 0DTE dynamics are captured only through end-of-day open interest residuals. Whether the structural patterns we detect transmit to overnight global markets---and whether such transmission would strengthen or weaken the regime signal---remains an open question requiring multi-venue data sources beyond the scope of this study.
